% ============================================================================
%  Section 1 -- Introduction
%
%  The marks in this section demonstrate the Mark API. Replace them with your
%  own content.
% ============================================================================
\section{Introduction}
\label{sec:introduction}

Inexpensive digital temperature sensors make dense monitoring networks
affordable, yet their factory specifications rarely match the accuracy that a
measurement study needs. Each unit deviates from the true temperature by an
offset and a gain error, and both terms drift over long
deployments~\cite{example2024drift}.
A useful correction must therefore estimate \MarkEdit{edits}{the offset}{both terms}
for every sensor.

Calibration against a reference instrument maps raw readings to values with a
stated uncertainty~\cite{taylor1997errors}. We follow the conventions of the
Guide to the Expression of Uncertainty in Measurement~\cite{jcgm2008gum}.
\Mark{changes}{The pipeline runs unattended once the sensors are in the bath, so a single operator can calibrate a whole batch in one session.}

\Mark{missing}{TODO: Add a sentence on prior calibration studies and cite them.}

\subsection{\Mark{changes}[heading]{Contributions}}

\begin{MarkEnv}{changes}[text]
We describe a two-parameter linear calibration model, a least-squares fit for its
two parameters, and a simple residual check that flags sensors whose response
is not linear over the calibrated range.
\end{MarkEnv}

\Cref{sec:method} presents the calibration model and the fitting procedure.
